% Fragmento para el informe existente. Requiere graphicx.
\section{Perspectiva del jugador (GM-1)}
Luma, guardiana de la llanura, descubre que la nave nodriza está contaminando la lluvia. El tarro del juego original (\texttt{bucket.png}) representa el recolector que lleva el protagonista. Las gotas azules (\texttt{drop.png}) representan agua pura y las contaminadas (\texttt{dropBad.png}) representan el peligro que introdujo la nave.

En la llanura aparece el primer módulo, el Pulso. En el sur helado, el protagonista puede recuperar una espada que ralentiza a la nave. En el oasis de Atacama puede encontrar el Láser Solar, que sigue causando daño después del impacto. Durante el Ascenso, el rayo del OVNI puede arrastrar una caja de hidro-plasma que dispara tres chorros cuando la barra está llena. La aparición depende del azar y el agua recogida alimenta cualquiera de estas armas. El protagonista puede completar la campaña con su chorro inicial.

El combate tiene tres fases. Primero la nave lanza ácido mientras se mueve horizontalmente. Al 65\,\% de vida, se enfurece y cambia su trayectoria y cadencia. Al 35\,\%, despliega dos drones y levanta un escudo. El jugador debe derribar a ambos drones mientras la nave se repara. Después del último dron, la nave vuelve a recibir daño. El aspecto Alien se desbloquea al ganar el encuentro oculto.

\section{Análisis de software (GM-2)}
El juego original movía un tarro y organizaba la información de la lluvia en arreglos separados. Eso obligaba a mantener sincronizados los índices de las gotas. Primero se introdujo una lista de objetos \texttt{Gota} y se delegó la detección de contacto y el efecto en cada variante. Luego \texttt{Partida} pasó a coordinar los niveles, las zonas y el combate. Así se pudo agregar un objeto especial sin replicar el ciclo de actualización de la lluvia.

Las armas terrestres se sortean con un 4\,\% por aparición normal, fuera del tramo de monedas. El Ascenso ofrece un único sorteo con 65\,\% de probabilidad. El último módulo recogido queda equipado. \texttt{ChorroAgua} conserva el arma usada al disparar, de modo que un cambio posterior no modifica un disparo en vuelo. Durante el combate siguen cayendo gotas de agua, escudos y botiquines. El escudo absorbe un golpe y el botiquín respeta el máximo de vidas de la dificultad.

\section{Diagramas de clases (GM-3)}
El primer diagrama muestra el conjunto de tipos de \texttt{core} y la clase base de LibGDX. Las dos figuras siguientes amplían las relaciones del dominio y de los sistemas que coordinan la interfaz, el combate y la persistencia.

\begin{figure}[!htbp]
\centering
\includegraphics[width=\linewidth,height=.72\textheight,keepaspectratio]{images/00-uml-completo.png}
\caption{UML maestro de GameLluvia.}
\end{figure}
\clearpage
\begin{figure}[!htbp]
\centering
\includegraphics[width=\linewidth,height=.72\textheight,keepaspectratio]{images/01-dominio.png}
\caption{Partida, controles, lluvia y reglas del dominio.}
\end{figure}
\clearpage
\begin{figure}[!htbp]
\centering
\includegraphics[width=\linewidth,height=.72\textheight,keepaspectratio]{images/06-sistemas.png}
\caption{Presentación, combate, Versus y perfil.}
\end{figure}
\clearpage

\section{Patrones de diseño (GM-6 a GM-9)}
\subsection{Singleton}
\textbf{Problema.} El menú y la partida necesitan consultar el mismo ajuste de sonido y el mismo perfil.
\textbf{Solución.} \texttt{ConfiguracionJuego} usa un constructor privado y \texttt{getInstancia()} crea la instancia al solicitarla por primera vez.
\textbf{Aplicación.} \texttt{GameLluvia.create()} obtiene esa configuración. El menú modifica el sonido y la partida consulta el ajuste para reproducir música y efectos.
\textbf{Ventaja y utilidad.} El cambio se refleja entre pantallas y el perfil se inicializa desde un punto común.

\subsection{Strategy dinámico}
\textbf{Problema.} El dispositivo de entrada, la superficie, el arma equipada y el comportamiento del jefe cambian durante el juego.
\textbf{Solución.} \texttt{Personaje} delega la entrada en \texttt{ControlMovimiento}, el desplazamiento en \texttt{MovimientoTerreno} y conserva una variante de \texttt{ArmaBioma}. Cada valor de \texttt{JefeOvni.Fase} define su trayectoria y acción.
\textbf{Aplicación.} T o el botón del HUD llaman a \texttt{setControlMovimiento()}. Al cambiar de tramo, \texttt{Partida.actualizarMovimientoTerreno()} llama a \texttt{setMovimientoTerreno()}. Una caja llama a \texttt{equiparArma()} mediante \texttt{atrapar()}. Al disparar, \texttt{JefeOvni.disparar()} usa el arma equipada para crear los chorros. \texttt{recibirDaño()} cambia la fase del jefe y \texttt{actualizar()} delega en la fase activa.
\textbf{Ventaja y utilidad.} Los cambios ocurren con los mismos objetos y mantienen su progreso. Los umbrales y las transiciones pertenecen al jefe, mientras que las variantes de trayectoria y ataque quedan en su fase. La curación se detiene al caer el último dron.

\subsection{Template Method}
\textbf{Problema.} Todas las gotas deben moverse, descartarse al salir y resolver el contacto una sola vez.
\textbf{Solución.} \texttt{Gota.actualizar()} es final y fija el orden. Llama a \texttt{calcularMovimiento()}, \texttt{verificarColision()} y \texttt{aplicarEfecto()}.
\textbf{Aplicación.} El agua cae a velocidad normal, usa el área completa y recarga al jugador. El ácido cae 15\,\% más rápido y reduce el área de contacto. \texttt{ObjetoEspecial} ajusta la velocidad según su tipo, usa un margen propio y delega el efecto a \texttt{Personaje.atrapar()}. Así se aplica tanto una mejora como la instalación de un arma.
\textbf{Ventaja y utilidad.} \texttt{Lluvia.mover()} actualiza la misma lista de gotas y elimina las terminadas. Cada subtipo personaliza varias etapas y el ciclo de vida permanece en la clase base.

\subsection{Builder}
\textbf{Problema.} Cada dificultad calcula varios parámetros por nivel y sus valores cambian a lo largo de la campaña.
\textbf{Solución.} \texttt{Nivel.Builder} permite nombrar cada ajuste antes de construir un \texttt{Nivel}.
\textbf{Aplicación.} \texttt{Dificultad.crearNivel()} configura el objetivo, la velocidad, el intervalo y la probabilidad de peligro. \texttt{Partida} lo invoca al cambiar de nivel.
\textbf{Ventaja y utilidad.} La configuración se puede leer en el orden en que se define, y la lluvia recibe un nivel construido con sus parámetros explícitos.


\section{Hidratación y carga del arma}
Tras cuatro segundos sin recoger agua, el Guardián empieza a perder velocidad. A los doce segundos recibe daño y luego vuelve a recibirlo cada cuatro segundos, hasta poder perder la partida. La pausa y el tramo de monedas detienen el desgaste. Una gota azul restablece el movimiento. La variante doble, identificada con x2, aporta dos cargas de munición y dos puntos a la barra del arma equipada.

La barra llega a seis con agua recogida después de equipar el arma. El siguiente disparo activa automáticamente el poder especial, consume la barra y gasta una carga de munición. La espada prolonga la ralentización, el láser aplica más pulsos de calor y el cañón lanza tres chorros. El jefe tiene 72, 90 o 108 puntos de vida según la dificultad.

\subsection{Evidencia de la recorrección}
\texttt{Personaje.actualizar()} selecciona la estrategia \texttt{Hidratacion.FATIGADA} y \texttt{recogerAgua()} restaura \texttt{NORMAL}. La estrategia de terreno y el control de entrada siguen funcionando al mismo tiempo. \texttt{JefeOvni.disparar()} lee el arma equipada y consume la carga cuando está lista. El proyectil conserva el arma y el estado de carga. El método final \texttt{Gota.actualizar()} ejecuta movimiento, colisión y efecto, y después llama al hook \texttt{despuesDeAtrapar()}. \texttt{GotaBuena} lo usa para avisar que el agua doble eliminó la fatiga y dio dos cargas.
